\documentclass[12pt]{article}
\usepackage[utf8]{inputenc}
\usepackage[margin=2.5cm]{geometry}

\begin{document}

\begin{flushright}
8 June 2026
\end{flushright}

\vspace{8pt}

\noindent Dear Editor,

\vspace{8pt}

\noindent We submit our manuscript ``Dark energy as an information capacity deficit'' for consideration as an Article in {\it Nature Physics}.

\vspace{8pt}

\noindent{\bf What we have found.}

Dark energy can be derived from two information-theoretic axioms. Structure formation in the universe occupies a finite reservoir of vacuum quantum information modes. As stars and galaxies form, the fraction $q(t)$ of modes remaining free decreases, and the effective dark energy density $\rho_{\rm DE}(t)=\rho_\Lambda q(t)$ evolves. The resulting equation of state $w(z)+1\propto{\rm SFR}(z)/[H(z)\rho_*(z)^{n/(n+1)}]$ is set by the cosmic star formation history with a single parameter. For the natural value of that parameter ($n=1$, damping proportional to the information deficit), the DGF prediction $(w_0,w_a)=(-0.80,-0.51)$ lies inside the DESI DR2 $1\sigma$ observed contour. The preference over $\Lambda$CDM is $\Delta\chi^2=56.9$.

The peer-reviewed dark energy papers triggered by DESI DR2---including those in {\it Nature Astronomy}~\cite{desi2025natast}---have parameterized $w(z)$ phenomenologically and fitted it to data. This paper provides a microphysical origin.

\vspace{8pt}

\noindent{\bf Why this belongs in {\it Nature Physics}.}

\begin{enumerate}
\item {\bf It is a conceptual advance, not a parameter fit.} CPL, quintessence, and emergent dark energy all choose a functional form for $w(z)$ and adjust its parameters to match data. DGF derives the functional form from axioms about the information capacity of spacetime. The agreement with DESI is a prediction, not a fit.

\item {\bf It makes a falsifiable prediction that no other model can produce.} $w(z)+1$ is predicted to have a peak at $z\sim1$--$2$ because the cosmic star formation rate itself peaks there. The CPL parameterization is linear in $(1-a)$ and cannot produce a peak for any choice of $(w_0,w_a)$. DESI DR3 and Euclid will test this directly.

\item {\bf It connects cosmology to quantum information theory.} The same two axioms that predict $w(z)$ also yield quantum mechanics, the Einstein equations, and the black hole entropy area law (summarized in the Supplementary Information). The $q$-field has a proposed laboratory test via quantum Darwinism in superconducting qubits.

\item {\bf It arrives at an opportune moment.} The DESI DR2 results have generated intense interest. Multiple phenomenological studies have appeared. A paper that provides a physical origin for the observed signal, rather than another fit, addresses the central open question raised by DESI.
\end{enumerate}

\vspace{8pt}

\noindent{\bf Manuscript information.}

The main text is approximately 3,200 words with three display items. The Supplementary Information (eight sections) contains the full derivation of the telegraph equation, the $w(z)$ shape function, the error budget, a competitor comparison table, and the quantum Darwinism laboratory test proposal. All numerical code and data are available at the project repository.

\vspace{8pt}

\noindent{\bf Suggested reviewers.}

\begin{enumerate}
\item A cosmologist with expertise in dark energy theory and DESI data.
\item A quantum information theorist working on emergent spacetime.
\item An observational cosmologist specializing in BAO or supernova cosmology.
\end{enumerate}

\vspace{8pt}

\noindent{\bf Competing interests.} None.

\vspace{8pt}

\noindent This work is not under consideration elsewhere.

\vspace{14pt}

\noindent Sincerely,\\
\vspace{8pt}
\noindent The Authors

\vspace{14pt}

\begin{thebibliography}{9}
\bibitem{desi2025natast} DESI Collaboration.\ Dynamical dark energy in light of DESI DR2. {\it Nature Astron.} {\bf 9}, 1879 (2025).
\end{thebibliography}

\end{document}
